%=====================================================================
\chapter{IUB Compliance Quick Reference: MS IT and PhD IT}\label{app:compliance}
%=====================================================================

Methodology does not exist in a vacuum. A design that is impeccable in the
abstract can still fail you if it cannot be completed inside your degree's
time limit, if its similarity report crosses a threshold, or if the paper it
produces lands in a journal category that does not count. This appendix
gathers the rules that actually bind you into one place, so that the rest of
the book can talk about method without stopping to recite regulation.

Two documents govern, and they are \emph{different documents with different
numbers}. Read the column for your own degree.

\begin{regbox}[Why this course exists --- PhD Degree Regulations-2024, \S10(iv)]
Beyond the four specialisation courses, ``the PhD student will be required to
take the following two courses of 3 credit hours each: (a) Advanced Statistical
Analysis/Hypothetical Testing/Research Methodology. (b) Dissertation Writing
including PhD Proposal writing.''

\medskip
This handout is course (a). \chref{synopsis} and \chref{writing} overlap
deliberately with course (b).
\end{regbox}

\begin{alertbox}[This appendix is a map, not the territory]
Regulations are amended. Every figure here was transcribed from the 2024
editions approved by the Syndicate on 13 June 2024, and every clause number is
given so you can check it. When a number here disagrees with the current
document from the Directorate of Advanced Studies and Research, \textbf{the
Directorate's document wins} --- and please tell the course instructor so this
appendix can be corrected.
\end{alertbox}

\section{The Two Governing Documents}

\begin{center}
\small
\begin{tabular}{@{}L{3.3cm}L{5.4cm}L{5.4cm}@{}}
\toprule
 & \textbf{PhD IT} & \textbf{MS IT} \\
\midrule
Governing document
  & \emph{Doctor of Philosophy (PhD) Degree Regulations-2024}~\cite{iub2024phd}
  & \emph{MPhil/MS/LLM/MSc (Hons)/MBA Degree Regulations-2024}~\cite{iub2024mphil} \\
NQF level
  & Level 8
  & Level 7 \\
Written product
  & Dissertation
  & Thesis \\
Formatting standard
  & \multicolumn{2}{L{11.1cm}}{\emph{IUB Theses/Synopses Guidelines-2023}~\cite{iub2023theses}, for both tracks} \\
\bottomrule
\end{tabular}
\end{center}

\section{Programme Structure}

\begin{center}
\small
\begin{longtable}{@{}L{3.5cm}L{5.5cm}L{5.5cm}@{}}
\toprule
\textbf{Item} & \textbf{PhD IT} & \textbf{MS IT} \\
\midrule
\endfirsthead
\toprule
\textbf{Item} & \textbf{PhD IT} & \textbf{MS IT} \\
\midrule
\endhead
\bottomrule
\endfoot
Coursework
  & 18 CH (\S41)
  & Research-based: 24 CH. Course-based: 36 CH (\S7, \S15) \\
Research / thesis credit
  & 36 CH
  & 6 CH (research-based); none (course-based) \\
Total degree
  & 54 CH
  & 30 CH or 36 CH \\
Load per semester
  & 9 CH (\S41)
  & 12 CH in semesters 1--2 (\S15) \\
Semester duration
  & 18 weeks; 16--18 with 2 weeks of examinations (\S41, \S53)
  & 18 weeks, 2 reserved for examinations (\S15, \S17) \\
Regular semesters
  & 16 maximum
  & 4 \\
Minimum duration
  & 3 years / 6 semesters (\S16)
  & 2 years / 4 semesters (\S13) \\
Maximum duration
  & 8 years / 16 semesters (\S16)
  & 4 years; 5 only on force majeure (\S13) \\
Class strength
  & Minimum 5 to commence (\S10(vii))
  & Minimum 8, maximum 40 (\S19) \\
\end{longtable}
\end{center}

\section{This Course: Assessment and Examination}

The weights are \emph{identical} across the two tracks, and they are fixed by
regulation --- an instructor cannot reallocate them. Appendix~\ref{app:schedule}
maps the course's portfolio of work onto these components rather than adding
anything on top of them.

\begin{center}
\small
\begin{tabular}{@{}L{5.4cm}cL{7.2cm}@{}}
\toprule
\textbf{Component} & \textbf{Weight} & \textbf{Basis} \\
\midrule
Sessional                       & 20\% & Four sub-components below \\
\quad Participation / behaviour / group work & 5\% & \\
\quad Quiz (surprise, 2--5 quizzes, 5 questions each) & 5\% & \\
\quad Assignment (at least one)  & 5\% & Marked on information and references
                                         included, logical reasoning, and
                                         organisation of material \\
\quad Presentation / seminar     & 5\% & 10--20 minutes, individual or group \\
Mid-term examination            & 30\% & After 7--8 weeks of teaching \\
Final-term examination          & 50\% & Whole syllabus; 20--30\% weight to
                                         pre-midterm material \\
\bottomrule
\end{tabular}
\end{center}

\noindent
\reg{PhD Regs 2024}{\S51.1--51.2}; \reg{MPhil/MS Regs 2024}{\S36}.

\subsection*{Paper patterns}

The mid-term is the same in both tracks. The final differs only in the clock.

\begin{center}
\small
\begin{tabular}{@{}L{4.2cm}L{3.3cm}cL{2.2cm}L{2.2cm}@{}}
\toprule
\textbf{Paper} & \textbf{Question type} & \textbf{Marks} & \textbf{Time (PhD)} & \textbf{Time (MS)} \\
\midrule
Mid-term (30 marks) & Objective, 10 or 20 & 10 & 15 min & 15 min \\
                    & Short answer, 5 at 2 marks & 10 & 25 min & 25 min \\
                    & Essay, 1 & 10 & 35 min & 35 min \\
\addlinespace
                    & \textbf{Total} & \textbf{30} & \textbf{1 h 15} & \textbf{1 h 15} \\
\midrule
Final (50 marks)    & Objective, 20 at 1 mark & 20 & 30 min & 20 min \\
                    & Short answer, 7 at 2 marks & 14 & 35 min & 40 min \\
                    & Essay, 2--4 & 16 & 55 min & 60 min \\
\addlinespace
                    & \textbf{Total} & \textbf{50} & \textbf{2 h} & \textbf{2 h} \\
\bottomrule
\end{tabular}
\end{center}

\begin{regbox}[Attendance --- PhD \S50, MPhil/MS \S35]
80\% attendance is required \emph{in each component} to sit the final
examination. Below 80\% is treated as failure in that course, with all the
consequences that has for good standing. Between 70\% and 80\% may be made up
by extra classes or assignments arranged by the teacher, but only for
unavoidable circumstances.
\end{regbox}

\subsection*{Grading --- where the two tracks genuinely diverge}

This is the difference students most often get wrong. The PhD scale is harder
in every respect: a higher pass mark, a higher CGPA floor, and a ``C'' that is
already disqualifying.

\begin{center}
\small
\begin{tabular}{@{}L{4.4cm}L{5.3cm}L{5.3cm}@{}}
\toprule
 & \textbf{PhD IT} (\S55, \S56) & \textbf{MS IT} (\S38, \S39) \\
\midrule
Pass percentage      & 65\%  & 60\% \\
Failing grade        & D, at 64\% or below & D, at 59\% or below \\
``C'' band           & 65--69\%, remarked \emph{Poor} & 60--69\%, remarked \emph{Satisfactory} \\
CGPA for the degree  & 3.00, with \textbf{no C grade} & 2.70, with \textbf{no D grade} \\
Probation after sem 1 & SGPA 2.70 & SGPA 2.30 \\
Dropped below        & CGPA 3.00 in later semesters & CGPA 2.70 in later semesters \\
Per-course floor     & GPA 2.75 in \emph{each} course before the comprehensive examination (\S11) & --- \\
\bottomrule
\end{tabular}
\end{center}

\begin{pitfall}[Reading the grade table backwards]
A PhD scholar who scores 67\% in this course has passed the course and
simultaneously acquired a ``C'' that disqualifies them from the degree until
it is repaired by repeating the course. The regulation permits repeating a
course to clear a C or to improve a B (\S56(v)). Notice also that fractional
marks round \emph{up} to the next whole figure: 64.10 becomes 65.00, which for
a PhD scholar is the difference between D and C.
\end{pitfall}

\section{Milestones After the Coursework}

\begin{center}
\small
\begin{longtable}{@{}L{3.5cm}L{5.5cm}L{5.5cm}@{}}
\toprule
\textbf{Milestone} & \textbf{PhD IT} & \textbf{MS IT} \\
\midrule
\endfirsthead
\toprule
\textbf{Milestone} & \textbf{PhD IT} & \textbf{MS IT} \\
\midrule
\endhead
\bottomrule
\endfoot
Supervisor
  & Consent at admission; appointed by the Board of Studies (\S21)
  & Assigned by end of semester 2 (\S8) \\
Comprehensive examination
  & Required. Within 60 days of completing coursework, at CGPA 3.0. Written
    and oral, 65\% to pass each. One retake within 30 days, then dropped (\S11)
  & \textbf{None}, except in the Institute of Business, Management and
    Administrative Sciences and the Department of Law (\S12) \\
Synopsis submission
  & Not later than completion of the 6th semester (\S13(iii))
  & To the supervisor within 8 weeks of the start of semester 3; DRC requires
    CGPA 2.70 first (\S8, \S10(i)) \\
Title approval
  & Board of Studies on DRC recommendation
  & Fixed by end of semester 3 (\S47(iii)) \\
Synopsis defence
  & Before the Departmental Research Committee, then BoS/BoF, then AS\&RB.
    Evaluated by at least one external referee from a national university,
    chosen from a list of 6 experts supplied by the supervisor (\S13)
  & Presented and defended before the Departmental Research Committee (\S8) \\
Publication required
  & \textbf{Yes.} One article in a W-category journal \emph{or} two in
    X-category, for Science disciplines. Scholar must be first author, the
    work must be relevant to the dissertation, and it must be published
    \emph{after} synopsis approval (\S15)
  & \textbf{No publication requirement} for the award of the degree \\
Examiners
  & Four (five with a co-supervisor): the supervisor, two foreign
    externals for evaluation, and one national external for the viva (\S33)
  & Two: the supervisor as internal, and one external chosen by the Dean from
    a panel of at least six (\S48) \\
Defence
  & Open defence, publicly announced by DASR, after positive external reports (\S33)
  & Viva voce in the presence of the external examiner; open defence preferred (\S48) \\
Verdicts
  & Pass / Pass with minor amendments / Deferred, re-defence within 180 days /
    Fail (\S33(vii))
  & Pass / Pass with minor amendments / Deferred for resubmission within
    six months / Fail (\S48(xi)) \\
\end{longtable}
\end{center}

\begin{conceptbox}[What this means for how you choose a method]
The two timetables reward different designs.

\medskip
An MS IT scholar has roughly two semesters of research time and no publication
gate. A design that yields a defensible answer in one pass --- a well-scoped
case study (\chref{casestudy}), a survey with an existing validated instrument
(\chref{survey}), a focused controlled experiment (\chref{experiments}) --- is
the sensible shape.

\medskip
A PhD scholar must publish before submitting, which means the research has to
produce a publishable increment \emph{part-way through}, not only at the end.
That favours a programme of linked studies: a systematic review that stands
alone as a paper (\chref{slr}), then the empirical or design-science work it
motivates (\chref{designscience}). Plan the first paper into the synopsis.
\end{conceptbox}

\section{Integrity Thresholds}

These are identical in both documents, and they are numeric, so there is no
room to argue about them after the fact.

\begin{regbox}[Similarity --- PhD \S32(III), MPhil/MS \S49(3)]
\begin{tight}
  \item Overall similarity index: \textbf{19\% or less}.
  \item From any \emph{single} source: \textbf{less than 5\%}.
  \item In the findings and conclusion section specifically: \textbf{not more
    than 9\%} --- ``the Similarity index should be considered very seriously in
    the section of findings and conclusion''.
\end{tight}
\medskip
A similarity test must be run before the thesis or dissertation goes to the
external experts. References, bibliography and tables of contents are removed
before submission for checking, because leaving them in inflates the index
(\S29 IV(iv)).
\end{regbox}

\begin{conceptbox}[Interpreting the report, not just the number]
The regulations are unusually careful here, and \chref{integrity} takes this
apart properly. A similarity index is not a plagiarism verdict: matches with
proverbs, universal truths, common phrases and proper nouns are expected, and
matches against your own earlier work may be ignored if the material is cited
(\S29 V). Conversely, an index of 8\% that is a single uncited paragraph from
one source is plagiarism, and an index of 18\% spread over correctly quoted and
cited material may not be. The instructor or supervisor delivers the verdict
after interpreting the report --- the software does not.
\end{conceptbox}

\section{Generative AI}

The IUB position is stricter than most international publisher policies and it
is internally in tension, so read both clauses rather than one.
\chref{genai} is devoted to working within it.

\begin{regbox}[Machine-generated text --- PhD \S32(I)(C)(vi), MPhil/MS \S49(C)(vi)]
Listed among activities that ``still count as academic cheating and plagiarism
if done without permission from the original artists/creators'':

\medskip
``(vi) Using ChatGPT and similar machine-generated text''.
\end{regbox}

\begin{regbox}[Permitted use --- PhD \S32(IV), MPhil/MS \S49(4)]
``Researchers may use ChatGPT, Artificial Intelligence (AI) and AI-assisted
technologies to understand basic phenomena of anything and should not replace
the key researcher tasks such as producing scientific insights, analyzing and
interpreting data or drawing scientific conclusions. The authors are
responsible and accountable for the contents of the work and should not rely
solely on AI-generated content.''
\end{regbox}

\noindent
Read together, the workable reading is: AI may help you \emph{understand};
it may not \emph{produce} the intellectual content, and it may not
generate text that then appears in your submission. You remain accountable
either way. \chref{genai} gives a permitted/prohibited table for this course
and a disclosure statement template, and Appendix~\ref{app:templates} repeats
the template in fill-in form.

\section{Research Ethics}

\begin{regbox}[Academic integrity --- PhD \S29(I)]
The regulations adopt the International Center for Academic Integrity's 2022
definition: academic integrity is ``commitment, even in the face of adversity,
to six fundamental values: honesty, trust, fairness, respect, responsibility,
and courage. From these values flow principles of behavior that enable academic
communities to translate ideals into action.''

\medskip
For publication and theses, COPE guidelines are to be followed
(\S29 IV(i))~\cite{cope2019corepractices}. Plagiarism cases are handled under
the HEC Anti-Plagiarism Policy~\cite{hec2023plagiarism}.
\end{regbox}

The dissertation itself is also assessed on methodological quality, and the
criteria are written into the regulations rather than left to the examiner.
\reg{PhD Regs 2024}{\S14.3} sets out separate checklists for qualitative and
quantitative work; \chref{qda} and \chref{validity} map the course's material
onto them directly, and Appendix~\ref{app:templates} turns them into a
pre-submission checklist.

\begin{conceptbox}[The four quality tests a PhD dissertation must satisfy]
\begin{tightnum}
  \item \textbf{Qualitative quality} (\S14.3A): does the research illuminate
    the subjective meaning, actions and context of those researched? Is the
    design responsive to real-life settings? Is the description detailed enough
    for a reader to interpret meaning and context?
  \item \textbf{Quantitative quality} (\S14.3B): reliability, validity,
    internal validity, external validity --- ecological and population --- and
    replicability.
  \item \textbf{Appropriateness of methods to aims} (\S14.3C): the methods
    justified against the objectives, and the results shown to support them.
  \item \textbf{Relevance to policy and practice} (\S14.3D): usefulness
    established, implications discussed, and the assessment ``must not be
    superficial and lacking substance''.
\end{tightnum}
\medskip
Note also \S14.1: the research area should correspond to community needs at
regional and local levels, comply with the national research agenda, and
preferably coincide with the Sustainable Development Goals.
\end{conceptbox}

\section{What This Appendix Does Not Cover}

Admission criteria, fee schedules, migration and credit transfer, semester
freezing, grievance procedure, sub-campus processing, and payment of foreign
evaluators are all in the source documents and are not reproduced here. If you
need them, go to the Directorate of Advanced Studies and Research rather than
to a second-hand summary --- including this one.
